% 316_Riemann_Zeta_En_ch.tex -- chapter content Doc. 316
% FFGFT and the Riemann zeta function: reach and limit

\section*{Preliminary: the answer in short}
\addcontentsline{toc}{section}{Preliminary: the answer in short}

\begin{keyresult}[Result]
The Riemann hypothesis is an \textbf{arithmetic} matter. The torus and
$\xi$ contribute nothing to its solution. The result of this document is
a \emph{negative} one -- three obvious geometric approaches are excluded
with reasons given.
\end{keyresult}

The blanket formulation would be too coarse in one respect, hence the
refinement.

\textbf{What remains purely a matter of numbers.} Where the zeros lie,
how they stand relative to one another, how one grasps them
individually -- that is carried by arithmetic, by the primes. Three
geometric approaches were tested and all are negative: no compact
geometry can carry the zeros as its spectrum (Section~\ref{sec:weyl}, at
theorem level), $\xi$ does not shift the critical line
(Section~\ref{sec:casimir}), and the $\xi$ ladder discretises nothing
(Section~\ref{sec:rand}, out-of-sample test negative).

\textbf{What is nevertheless not merely a matter of numbers.} Two
things. First, the Riemann zeta is a \emph{factor} in the spectrum of
the $\mathbb{Z}^4$ torus -- an identity, not a picture
(Section~\ref{sec:faktor}). The Riemann hypothesis is thereby a
statement about a geometric object of the framework; only one that the
geometry cannot answer. Second, the form of the zero density says
\emph{where} the question belongs: $N(T)$ grows like $T\ln T$, and no
Laplacian on any compact manifold can deliver that -- it requires a
dilation structure, hence the unrolled domain. The geometry indicates
the address without supplying the answer.

\textbf{The striking thing is precisely that both occur together:} the
numbers demonstrably sit in the torus spectrum, but the torus cannot
generate them.

\textbf{The yield.} Three dead ends are closed, with reasons rather than
conjecture. In addition, one of the author's own faulty constructions is
documented: a cell argument that initially looked load-bearing turned
out on elaboration to be a restatement of $\hbar = 1$
(Section~\ref{sec:zelle}) and was withdrawn.

\section{Starting point and claim}

\textbf{[Q]} The corpus mentions the Riemann hypothesis in exactly one
place: Doc.~024 notes, in an AI/network context, that the
\enquote{number space} has dimensions \enquote{modulated by the Riemann
hypothesis and $\zeta$ functions}. That is a conceptual marginal remark
without an epistemic marker and without derivation. A worked-out
connection does not exist in the corpus.

\textbf{[K]} At the same time Doc.~314 contains material that touches
the $\zeta$ function directly: Epstein zeta, Casimir point, heat kernel
on $\text{T}^4$. This document examines systematically what follows.

\textbf{Claim.} It proves nothing about the Riemann hypothesis and does
not claim to. Three things are clarified: which connections hold as
identities \textbf{[K]}, which path carries constructively \textbf{[B]},
and which obvious paths are \emph{demonstrably} not passable
\textbf{[X]}. The exclusion is the most valuable part -- it saves work
on dead ends.

\section{Two findings that hold as identities}
\label{sec:faktor}

\subsection{The Riemann zeta is a factor of the torus spectral zeta}

\textbf{[K]} Doc.~314 carries the closed form of the $\mathbb{Z}^4$
spectral zeta (verified threefold there; confirmed here against the
direct sum with analytic tail correction to $10^{-7}$ through
$10^{-12}$):
\[
  Z_{\mathbb{Z}^4}(s) \;=\; 8\,\bigl(1 - 4^{1-s}\bigr)\,
  \zeta(s)\,\zeta(s-1).
\]
This is not an analogy but an identity. The zeros of the torus spectral
zeta therefore have three sources:

\begin{center}
\begin{tabular}{ll}
\toprule
Source & Location of the zeros \\
\midrule
$\zeta(s) = 0$ & $\Re(s) = 1/2$ (non-trivial -- the RH) \\
$\zeta(s-1) = 0$ & $\Re(s) = 3/2$ (shifted by 1) \\
$1 - 4^{1-s} = 0$ & $s = 1 + 2\pi i k/\ln 4$ (lattice factor) \\
\bottomrule
\end{tabular}
\end{center}

On this reading the Riemann hypothesis is a statement about the spectrum
of the $\mathbb{Z}^4$ torus -- about a geometric object of the
framework, not merely about a number series.

\subsection{The Casimir point lies symmetric to the critical line}
\label{sec:casimir}

\textbf{[K]} Doc.~314 computes the $\zeta$-regularised vacuum energy
$E = \pi\,\zeta_M(-1/2)$ and finds the order reversal: D4 minimises the
Epstein zeta for $s > 0$, at $s = 0$ all lattices are equal, and at
$s = -1/2$ the $\mathbb{Z}^4$ torus lies lower ($-0.9321$ against
$-0.8686$).

The Casimir point $s = -1/2$ has, under the functional equation
$\zeta(s) = \chi(s)\zeta(1-s)$, the mirror point $s = 3/2$. The axis of
symmetry:
\[
  \frac{-1/2 + 3/2}{2} \;=\; \frac12
  \qquad\text{-- exactly the critical line.}
\]

The physically occupied interval $[-1/2,\,3/2]$ has the critical line as
its midpoint: consistency between physics and arithmetic.

\textbf{[X]} No proof follows from this -- the critical line is fixed by
the functional equation anyway. Nor can $\xi$ shift it: a perturbation
of order $\xi \sim 10^{-4}$ (even with the RG amplifier $100\xi$) does
not reach an axis fixed by a functional identity.

\section{The decisive obstruction: Weyl against Riemann-von Mangoldt}
\label{sec:weyl}

The obvious move would be to seek an operator whose eigenvalues are the
zeros (Hilbert-P\'olya programme), and to deform
$\text{T}^4/\mathbb{Z}_3$ suitably for that purpose. \textbf{This entire
branch is excluded -- at theorem level, not as intuition.}

\textbf{[Q]} \emph{Weyl's law.} For \emph{every} compact Riemannian
manifold of any dimension $d$, asymptotically
$N(\lambda) \sim C_d\,\text{Vol}\,\lambda^{d/2}$ -- a pure power.
Curvature, boundary and orbifold points change only lower orders, not
the leading term.

\textbf{[Q]} \emph{Riemann-von Mangoldt (proved).}
\[
  N(T) = \frac{T}{2\pi}\ln\frac{T}{2\pi} - \frac{T}{2\pi}
  + \frac78 + O(1/T).
\]

\textbf{[K]} The comparison via the local exponent
$\mathrm d\ln N/\mathrm d\ln T$:

\begin{center}
\begin{tabular}{rrr}
\toprule
$T$ & $N(T)$ & local exponent $d/2$ \\
\midrule
$10^{3}$ & $6.48\cdot10^{2}$ & $1.2457$ \\
$10^{5}$ & $1.38\cdot10^{5}$ & $1.1153$ \\
$10^{8}$ & $2.48\cdot10^{8}$ & $1.0642$ \\
$10^{12}$ & $3.95\cdot10^{12}$ & $1.0403$ \\
$10^{20}$ & --- & $1.0233$ \\
\bottomrule
\end{tabular}
\end{center}

The exponent is $1 + 1/\ln(T/2\pi)$ -- it falls monotonically towards 1
but is \emph{never constant}. Range over 17 decades: $0.22$. The best
power fit in the window $[10^3, 10^6]$ deviates systematically by up to
$11\,\%$.

\begin{keyresult}[Exclusion at theorem level]
\textbf{[X]} There is no compact Riemannian manifold -- of whatever
dimension, curvature, topology or orbifold structure -- whose Laplace
spectrum is the Riemann zeros. Reason: $T\ln T$ is not a power
$T^{d/2}$.
\end{keyresult}

\textbf{Consequence.} The branch \enquote{deform
$\text{T}^4/\mathbb{Z}_3$ suitably} is settled completely, not only in
the flat case. The extra $\ln T$ is the signature of a \textbf{scale
degree of freedom}: a dilation direction with logarithmically growing
volume. That is not a Laplacian on a manifold.

\textbf{Withdrawn intermediate diagnosis.} In a first pass it was argued
that $\log(p)$ geodesics require negative curvature. That is wrong: what
generates log lengths is a dilation structure (multiplicative group
$\mathbb{R}_+^*$), not curvature; the Berry-Keating operator $H = xp$ is
the generator of dilation, not of a hyperbolic flow. The Weyl
obstruction is the correct and considerably stronger statement, and it
renders the curvature question moot.

\section{The constructive link -- and why it does not carry}
\label{sec:zelle}

If the zero density requires a scale degree of freedom, the responsible
domain is the \textbf{unrolled} one (Doc.~315) -- multiplicative,
logarithmic, with dilation as its natural operation.

\subsection{The approach}

\textbf{[Q]} \emph{Berry-Keating (1999).} The semiclassical count of the
dilation generator $H = (xp+px)/2$ in a phase space with a minimal cell:
\[
  N(E) = \frac{1}{2\pi}\bigl[\,E\ln\bigl(E/(l_x l_p)\bigr)
  - E + l_x l_p\,\bigr].
\]
With $l_x l_p = 2\pi$ this becomes exactly the Riemann-von Mangoldt
formula (verified: difference constant $0.125 = 1 - 7/8$, a pure
boundary term, over $T = 15 \ldots 5000$).

The obvious chain would then be: Doc.~314 Ch.~D1 carries
$\lambda_4\cdot m = 2\pi$ as an exact position relation; in the unrolled
domain the quantity conjugate to scale is the dilation momentum; hence
FFGFT supplies the cell condition.

\subsection{Why the chain collapses}

\textbf{[X] Stage 1 -- the cell condition is empty of content.}
$l_x l_p = 2\pi\hbar$ is, in natural units, the Planck cell, and follows
from the canonical quantisation of \emph{every} conjugate pair. That
$\lambda_4 m = 2\pi$ has the same form is no independent finding -- the
de Broglie relation \emph{is} the statement $\hbar = 1$. Moreover the
count is logarithmically insensitive to the cell (a factor 2 wrong
$\Rightarrow$ $N$ changes by ${\sim}1\,\%$); a \enquote{hit} would not
have been a test in any case.

\textbf{[X] Stage 2 -- the actual content fails quantitatively.}
Berry-Keating requires something sharper than the cell: two
\emph{independent} cutoffs $x > l_x$ and $p > l_p$ whose product is the
Planck cell. The FFGFT double cutoff makes this testable:

\begin{center}
\begin{tabular}{lr}
\toprule
Quantity & Value \\
\midrule
UV: $L_0 = \xi\,\ell_P$ & $2.155\cdot10^{-39}$ m \\
IR: $L^* = 4\bar\lambda_e/\xi^{10}$ & $8.698\cdot10^{26}$ m \\
Ratio $L^*/L_0$ & $4.04\cdot10^{65}$ \\
required & $2\pi \approx 6.28$ \\
\midrule
Shortfall & \textbf{65 decades} \\
\bottomrule
\end{tabular}
\end{center}

\textbf{[X] Stage 3 -- a finite window against infinitely many zeros.}
$\ln(L^*/L_0) = 151.06$, that is $16.93$ $\xi$ ladder rungs; the
Berry-Keating count within the window gives ${\sim}1816$ modes. A finite
scale window carries finitely many modes; the zeros are countably
infinite. Even if stages~1 and~2 held, at most an initial segment would
result -- and an initial segment is not a spectral identity.

\emph{Marginal observation, explicitly not booked:} $16.93$ lies close
to 17. Both cutoffs carry their own uncertainties -- the prefactor~4 is
open per R73, $\ell_P$ depends on $G$. Without an error budget this is
not assessable under P35.

\subsection{What remains}

\textbf{To be withdrawn} is the formulation \enquote{the zero density
follows from the FFGFT cell condition}. It suggests a finding where a
triviality stands.

\textbf{[B] A structural statement remains, and it is not little:} that
the counting function has the form $(T/2\pi)\ln(T/2\pi)$ at all requires
a \textbf{dilation generator}. A Laplacian on a compact manifold cannot
deliver it (Section~\ref{sec:weyl}). The substantive statement is the
\emph{identification of the structure} -- the unrolled domain carries
the dilation -- not the numerical value $2\pi$.

\section{The missing boundary condition -- test negative}
\label{sec:rand}

$H = xp$ has \emph{continuous} spectrum. The counting function fixes the
mean density, not the individual zeros. What is missing is the boundary
condition that turns the continuum into a discrete spectrum.

\textbf{The natural candidate.} The $\xi$ ladder generates a discrete
subgroup $\Gamma_\xi = \{\xi^n\} < \mathbb{R}_+^*$; the quotient
$\mathbb{R}_+^*/\Gamma_\xi$ is a \emph{circle} of circumference
$L = \ln(1/\xi) = 8.9227$. Admissible dilation frequencies would then be
multiples of $2\pi/L = 0.7042$.

\subsection{Test discipline}

\textbf{[K]} Proximity arguments require three controls (P35): a null
distribution over random periods in the \emph{same} range; artefact
control, because periods larger than the data range automatically
produce small scores; and an out-of-sample check with look-elsewhere
correction.

The artefact control is not theoretical: a first run reported an
apparent signal at period $2\pi/|\ln(74/75)| = 468$. Since all zeros lie
$< 102$, all quotients round to 0 -- a pure artefact of a null
distribution that did not cover the range. Documented because exactly
this error type recurs with $\xi$ proximity arguments.

\subsection{Result}

\textbf{[K]} Candidates within the admissible range (null distribution:
mean $0.2502$):

\begin{center}
\begin{tabular}{lrrr}
\toprule
Candidate & Period & Score & $p$ value \\
\midrule
$2\pi/\ln(1/\xi)$ & $0.7042$ & $0.2812$ & $0.92$ \\
$2\pi/\ln 75$ & $1.4553$ & $0.2043$ & $\mathbf{0.005}$ \\
$\ln(1/\xi)$ & $8.9227$ & $0.2891$ & $0.94$ \\
$2\pi$ & $6.2832$ & $0.2366$ & $0.18$ \\
\bottomrule
\end{tabular}
\end{center}

The borderline candidate $2\pi/\ln 75$ was checked out of sample (zeros
31--50): score $0.2686$, $p = 0.81$; look-elsewhere corrected
$p = 1.0$.

\begin{keyresult}[Boundary condition]
\textbf{[X]} The $\xi$ ladder does not supply the missing boundary
condition. No $\xi$-derived period candidate survives the out-of-sample
check.
\end{keyresult}

\subsection{Why this was structurally to be expected}

\textbf{[K]} A single-period compactification produces an
\emph{equidistant} lattice. The zeros, however, show GUE statistics with
level repulsion (Montgomery-Odlyzko): mean normalised spacing $0.9815$
(should be ${\sim}1$), fraction of spacings $< 0.3$ equal to
\textbf{$0.000$} -- against ${\sim}0.04$ for GUE and ${\sim}0.26$ for a
lattice-like spectrum. A circle cannot deliver GUE in principle.

The missing boundary condition is carried by arithmetic (the primes),
not by geometry.

\section{Balance}

\begin{center}
\begin{tabular}{ll}
\toprule
Statement & Status \\
\midrule
$\zeta$ is a factor of the torus spectral zeta & \textbf{[K]} verified \\
Casimir interval has critical line as axis & \textbf{[K]} \\
Zero density requires dilation generator & \textbf{[B]} structure \\
Cell condition as an independent finding & \textbf{[X]} only $\hbar = 1$ \\
Double cutoff as BK truncation & \textbf{[X]} 65 decades \\
Compact manifold with zero spectrum & \textbf{[X]} Weyl \\
$\xi$ ladder discretises the zeros & \textbf{[X]} test negative \\
$\xi$ shifts the critical line & \textbf{[X]} functionally fixed \\
FFGFT solves the Riemann hypothesis & \textbf{[X]} \\
\bottomrule
\end{tabular}
\end{center}

\textbf{The most honest sentence.} Of the three connections examined,
two hold as identities \textbf{[K]} -- but they are consistency, not a
statement about the location of the zeros. The third, the constructive
one, collapsed on elaboration: what looked like a derivation is
$\hbar = 1$. What remains is a structural statement: the zero density
requires a dilation generator, and that lives in the unrolled domain.

\textbf{What a viable approach would have to achieve:} a non-trivial
boundary condition on the unrolled scale direction that (a) carries
infinitely many modes -- hence is not a finite cutoff window -- and
(b) produces GUE statistics rather than equidistance. Both requirements
point in the same direction: the structure must be arithmetic, not
geometric.

\section{Open items}

\begin{itemize}
  \item \textbf{Z-1:} Does the $\xi$ ladder carry information about the
    Euler product? $\ln(1/\xi) = 2\ln 2 + \ln 3 + 4\ln 5$ decomposes
    exactly into the prime factors of 30000. Whether that is more than
    the trivial prime factorisation: open, no test formulated so far.
  \item \textbf{Z-2:} The lattice factor $(1 - 4^{1-s})$ yields zeros at
    $s = 1 + 2\pi i k/\ln 4$. Does this ladder have a corpus meaning
    (D4? the factor 4 from $L^* = 4\bar\lambda_e/\xi^{10}$, R73)? Not
    yet examined.
  \item \textbf{Z-3:} The heat-kernel finding of Doc.~314
    (D4/$\mathbb{Z}^4$ difference exponentially small,
    $1.8\cdot10^{-13}$ at $t = 0.1$) has a counterpart in $\zeta$ theory
    (theta-series fine structure). Whether a test surface follows:
    open.
  \item \textbf{Not recommended:} further search for $\xi$ periods in
    the zeros. Section~\ref{sec:rand} excludes this structurally via the
    GUE argument, not merely numerically.
\end{itemize}

\section{Verification}

All numbers are reproducible by five scripts (standard library only,
target values as assertions, seed 20780458):

\begin{itemize}
  \item \texttt{z1\_weyl\_obstruktion.py} -- Weyl's law against
    Riemann-von Mangoldt, local exponent, power fit.
  \item \texttt{z2\_zellbedingung.py} -- Berry-Keating count,
    equivalence to the Riemann form, cross-check against Weyl.
  \item \texttt{z3\_randbedingung\_test.py} -- $\xi$ ladder test with
    artefact control, null distribution, out-of-sample, GUE control.
  \item \texttt{z4\_epstein\_casimir.py} -- factorisation against the
    direct sum, Casimir symmetry point.
  \item \texttt{z5\_zellbedingung\_kritik.py} -- collapse of the cell
    argument in three stages.
\end{itemize}

Inputs: $\xi = 4/30000$ exact; $\ell_P = 1.616255\cdot10^{-35}$~m,
$\bar\lambda_e = 3.8615926796\cdot10^{-13}$~m (CODATA~2022); zero values
(first 50) from the literature.
